\documentclass{beamer}
\usepackage[utf8]{inputenc}

\usetheme{Madrid}
\usecolortheme{default}
\usepackage{amsmath,amssymb,amsfonts,amsthm}
\usepackage{txfonts}
\usepackage{tkz-euclide}
\usepackage{listings}
\usepackage{adjustbox}
\usepackage{array}
\usepackage{tabularx}
\usepackage{gvv}
\usepackage{lmodern}
\usepackage{circuitikz}
\usepackage{lmodern}
\usepackage{multicol}
\usepackage{tikz}
\usepackage{graphicx}
\usepackage{mathtools}
\setbeamertemplate{page number in head/foot}[totalframumber]

\usepackage{tcolorbox}
\tcbuselibrary{minted,breakable,xparse,skins}



\definecolor{bg}{gray}{0.95}
\DeclareTCBListing{mintedbox}{O{}m!O{}}{%
  breakable=true,
  listing engine=minted,
  listing only,
  minted language=#2,
  minted style=default,
  minted options={%
    linenos,
    gobble=0,
    breaklines=true,
    breakafter=,,
    fontsize=\small,
    numbersep=8pt,
    #1},
  boxsep=0pt,
  left skip=0pt,
  right skip=0pt,
  left=25pt,
  right=0pt,
  top=3pt,
  bottom=3pt,
  arc=5pt,
  leftrule=0pt,
  rightrule=0pt,
  bottomrule=2pt,
  toprule=2pt,
  colback=bg,
  colframe=orange!70,
  enhanced,
  overlay={%
    \begin{tcbclipinterior}
    \fill[orange!20!white] (frame.south west) rectangle ([xshift=20pt]frame.north west);
    \end{tcbclipinterior}},
  #3,
}
\lstset{
    language=C,
    basicstyle=\ttfamily\small,
    keywordstyle=\color{blue},
    stringstyle=\color{orange},
    commentstyle=\color{green!60!black},
    numbers=left,
    numberstyle=\tiny\color{gray},
    breaklines=true,
    showstringspaces=false,
}
%------------------------------------------------------------
%This block of code defines the information to appear in the
%Title page
\title %optional
{12.536}
\date{October 10,2025}
%\subtitle{A short story}

\author % (optional)
{EE25BTECH11002 - Achat Parth Kalpesh}



\begin{document}


\frame{\titlepage}

\begin{frame}{Question}
The rank of the matrix 
$
\myvec{
1 & 1 & 0 & -2 \\
2 & 0 & 2 & 2 \\
4 & 1 & 3 & 1 \\
}$
is

\begin{multicols}{4}
\begin{enumerate}
    \item 1
    \item 2
    \item 3
    \item 4
\end{enumerate}
\end{multicols}
\end{frame}

\begin{frame}{Solution}
We perform row operations to reduce it to row echelon form:


\begin{align}
\myvec{
1 & 1 & 0 & -2 \\
2 & 0 & 2 & 2 \\
4 & 1 & 3 & 1
}
&\xleftrightarrow[]{R_2 \leftarrow R_2 - 2R_1}
\myvec{
1 & 1 & 0 & -2 \\
0 & -2 & 2 & 6 \\
4 & 1 & 3 & 1
} \\
\myvec{
1 & 1 & 0 & -2 \\
0 & -2 & 2 & 6 \\
4 & 1 & 3 & 1
}&\xleftrightarrow[]{R_3 \leftarrow R_3 - 4R_1}
\myvec{
1 & 1 & 0 & -2 \\
0 & -2 & 2 & 6 \\
0 & -3 & 3 & 9
} \\
\myvec{
1 & 1 & 0 & -2 \\
0 & -2 & 2 & 6 \\
0 & -3 & 3 & 9
}&\xleftrightarrow[]{R_2 \leftarrow -\frac{1}{2} R_2}
\myvec{
1 & 1 & 0 & -2 \\
0 & 1 & -1 & -3 \\
0 & -3 & 3 & 9
} \\
\myvec{
1 & 1 & 0 & -2 \\
0 & 1 & -1 & -3 \\
0 & -3 & 3 & 9
}&\xleftrightarrow[]{R_1 \leftarrow R_1 - R_2}
\myvec{
1 & 0 & 1 & 1 \\
0 & 1 & -1 & -3 \\
0 & -3 & 3 & 9
}
\end{align}
\end{frame}



\begin{frame}{Solution}
\begin{align}
    \myvec{
1 & 0 & 1 & 1 \\
0 & 1 & -1 & -3 \\
0 & -3 & 3 & 9
}&\xleftrightarrow[]{R_3 \leftarrow R_3 + 3 R_2}
\myvec{
1 & 0 & 1 & 1 \\
0 & 1 & -1 & -3 \\
0 & 0 & 0 & 0
}
\end{align}
Now the matrix is in reduced row echelon form \\
The number of non-zero rows = 2

\begin{align}
    \text{rank}\brak{A} = 2
\end{align}
\end{frame}

\end{document}